% Half-title, title and copyright pages.
%
% \bookversion is the release the PDF was built for: `make VERSION=v1.1.0` passes it in. Without it
% the title page carries only the date of the build.
\providecommand{\bookversion}{}
\newcommand{\bookdate}{\number\day\space\ifcase\month\or January\or February\or March\or April\or
  May\or June\or July\or August\or September\or October\or November\or December\fi\space
  \number\year}
\newcommand{\bookedition}{\ifx\bookversion\empty\else\bookversion\,·\,\fi\bookdate}

\thispagestyle{empty}
\vspace*{2.2in}
\begin{flushright}
  {\Large\bfseries Digital Design with VHDL}
\end{flushright}
\cleardoublepage

\thispagestyle{empty}
\vspace*{1.1in}
\begin{flushleft}
  {\color{accent}\rule{0.9in}{2.5pt}}\par\vspace{1.4ex}
  {\fontsize{32}{38}\selectfont\bfseries Digital Design\\[0.15ex] with VHDL\par}
  \vspace{2.2ex}
  {\Large\itshape From a gate network on paper\\to a design running on a real FPGA\par}
  \vspace{1.2in}
  {\large Erik Pihl\par}
  \vfill
  {\small\scshape\lsstyle qrtech academy\par}
  \vspace{0.6ex}
  {\small\color{muted} \bookedition\par}
\end{flushleft}
\newpage

\thispagestyle{empty}
\vspace*{\fill}
{\small
\noindent\textit{Digital Design with VHDL}\\
Copyright \textcopyright\ 2026 QRTECH Academy.

\medskip
\noindent The text and figures of this book, and of the course it is typeset from, are licensed
under the Creative Commons Attribution-NonCommercial-ShareAlike 4.0 International License
(CC BY-NC-SA 4.0). You may copy and redistribute them, and adapt them, for any non-commercial
purpose, provided you give appropriate credit and share what you make under the same license. To
view a copy of the license, visit \url{https://creativecommons.org/licenses/by-nc-sa/4.0/}.

\medskip
\noindent The code in this book may also be used under the MIT License, which is the license of all
the code in the companion repository. The repository holds the lecture material, the example
designs, the testbenches and the self-assessment papers:

\medskip
\noindent\url{https://github.com/qrtech-academy/digital-design-vhdl}

\medskip
\noindent The VHDL in this book is VHDL-93 throughout, simulated with GHDL and synthesized with
Quartus Prime Lite for the Terasic DE0-CV board and its Cyclone~V FPGA. Quartus, Cyclone,
DE0-CV and CircuitVerse are trademarks of their respective owners.

\medskip
\noindent Set in TeX Gyre Pagella and DejaVu Sans Mono with Lua\LaTeX. The figures are drawn from
code, in the repository's \file{diagrams/}, except the circuits exported from CircuitVerse.

\medskip
\noindent\textcolor{muted}{Typeset from the eight lectures of the course, their exercises, and the
two self-assessment papers. This edition: \bookedition.}
\par}
\cleardoublepage
